\documentclass[11pt]{article}
\usepackage[margin=1in]{geometry}
\usepackage{amsmath,amssymb,booktabs,listings,xcolor}

\title{UR5: Forward Kinematics}
\author{RBE 501}
\date{}

\lstset{
    language=Matlab,
    backgroundcolor=\color{black!5},
    basicstyle=\footnotesize\ttfamily,
    keywordstyle=\color{blue},
    commentstyle=\color{green!60!black},
    stringstyle=\color{purple},
    breaklines=true,
    showstringspaces=false
}

\begin{document}
\maketitle

%=====================================================================
\section{Product of Exponentials (POE) Forward Kinematics}

In the space-frame Product of Exponentials (POE) approach, the continuous configuration state of the end-effector relative to the base frame is described by:
\begin{equation}
  T(\mathbf{\theta}) = e^{[\mathcal{S}_1]\theta_1} e^{[\mathcal{S}_2]\theta_2} e^{[\mathcal{S}_3]\theta_3} e^{[\mathcal{S}_4]\theta_4} e^{[\mathcal{S}_5]\theta_5} e^{[\mathcal{S}_6]\theta_6} M
\end{equation}

\paragraph{Home Position Matrix ($M$).}
When all joint angles are zero ($\mathbf{\theta}=\mathbf{0}$), the manipulator is in a fully extended vertical position. The initial end-effector pose matrix $M \in SE(3)$ is given by:
\begin{equation}
  \boxed{\;
  M = \begin{bmatrix} 1 & 0 & 0 & 0 \\ 0 & 1 & 0 & d_4+d_6 \\ 0 & 0 & 1 & d_1+a_2+a_3 \\ 0 & 0 & 0 & 1 \end{bmatrix}
  \;}
\end{equation}

\paragraph{Joint 1.}
The rotation axis points along the positive $Z$-axis, passing through the origin.
\begin{equation}
  \mathbf{\omega}_1 = \begin{bmatrix} 0 \\ 0 \\ 1 \end{bmatrix}, \quad \mathbf{q}_1 = \begin{bmatrix} 0 \\ 0 \\ 0 \end{bmatrix} \quad \Rightarrow \quad \boxed{\;\mathcal{S}_1 = \begin{bmatrix} 0 & 0 & 1 & 0 & 0 & 0 \end{bmatrix}^T\;}
\end{equation}

\paragraph{Joint 2.}
The axis is parallel to the $Y$-axis, offset by $d_4$ laterally and $d_1$ vertically.
\begin{equation}
  \mathbf{\omega}_2 = \begin{bmatrix} 0 \\ 1 \\ 0 \end{bmatrix}, \quad \mathbf{q}_2 = \begin{bmatrix} 0 \\ d_4 \\ d_1 \end{bmatrix} \quad \Rightarrow \quad \boxed{\;\mathcal{S}_2 = \begin{bmatrix} 0 & 1 & 0 & d_1 & 0 & 0 \end{bmatrix}^T\;}
\end{equation}

\paragraph{Joint 3.}
The axis is parallel to the $Y$-axis, located at the upper arm height.
\begin{equation}
  \mathbf{\omega}_3 = \begin{bmatrix} 0 \\ 1 \\ 0 \end{bmatrix}, \quad \mathbf{q}_3 = \begin{bmatrix} 0 \\ d_4 \\ d_1+a_2 \end{bmatrix} \quad \Rightarrow \quad \boxed{\;\mathcal{S}_3 = \begin{bmatrix} 0 & 1 & 0 & d_1+a_2 & 0 & 0 \end{bmatrix}^T\;}
\end{equation}

\paragraph{Joint 4.}
The axis is parallel to the $Y$-axis, located at the forearm length height.
\begin{equation}
  \mathbf{\omega}_4 = \begin{bmatrix} 0 \\ 1 \\ 0 \end{bmatrix}, \quad \mathbf{q}_4 = \begin{bmatrix} 0 \\ 0 \\ d_1+a_2+a_3 \end{bmatrix} \quad \Rightarrow \quad \boxed{\;\mathcal{S}_4 = \begin{bmatrix} 0 & 1 & 0 & d_1+a_2+a_3 & 0 & 0 \end{bmatrix}^T\;}
\end{equation}

\paragraph{Joint 5.}
The axis is parallel to the $Z$-axis at home position, sharing the same vertical height as Joint 4.
\begin{equation}
  \mathbf{\omega}_5 = \begin{bmatrix} 0 \\ 0 \\ 1 \end{bmatrix}, \quad \mathbf{q}_5 = \begin{bmatrix} 0 \\ d_4 \\ d_1+a_2+a_3 \end{bmatrix} \quad \Rightarrow \quad \boxed{\;\mathcal{S}_5 = \begin{bmatrix} 0 & 0 & 1 & d_4 & 0 & 0 \end{bmatrix}^T\;}
\end{equation}

\paragraph{Joint 6.}
The terminal rotation axis is parallel to the $Y$-axis.
\begin{equation}
  \mathbf{\omega}_6 = \begin{bmatrix} 0 \\ 1 \\ 0 \end{bmatrix}, \quad \mathbf{q}_6 = \begin{bmatrix} 0 \\ d_4 \\ d_1+a_2+a_3+d_5 \end{bmatrix} \quad \Rightarrow \quad \boxed{\;\mathcal{S}_6 = \begin{bmatrix} 0 & 1 & 0 & d_1+a_2+a_3+d_5 & 0 & 0 \end{bmatrix}^T\;}
\end{equation}

%=====================================================================
\section{Geometric Jacobian Formulation}

The $6 \times 6$ Geometric Jacobian matrix $J(\mathbf{q})$ maps the joint velocities directly to the end-effector spatial velocities. It is constructed column by column based on the current axis orientation vectors $\mathbf{z}_{i-1}$ and link frame origins $\mathbf{o}_{i-1}$ extracted from the forward kinematics steps:
\begin{align}
  T_{01} &= e^{[\mathcal{S}_1]\theta_1} M, \quad T_{02} = T_{01}e^{[\mathcal{S}_2]\theta_2}, \quad \dots, \quad T_{06} = T(\mathbf{\theta})
\end{align}
Let $\mathbf{o}_{end} = T_{06}(1:3, 4)$ represent the instantaneous position of the end-effector tip.

\paragraph{Jacobian Columns Mapping.}
\begin{align}
  \text{Column 1 (Joint 1):} \quad & J_1 = \begin{bmatrix} \mathbf{z}_0 \times (\mathbf{o}_{end} - \mathbf{o}_0) \\ \mathbf{z}_0 \end{bmatrix}, \quad \text{where } \mathbf{z}_0 = \begin{bmatrix}0\\0\\1\end{bmatrix}, \, \mathbf{o}_0 = \begin{bmatrix}0\\0\\0\end{bmatrix} \\
  \text{Column 2 (Joint 2):} \quad & J_2 = \begin{bmatrix} \mathbf{z}_1 \times (\mathbf{o}_{end} - \mathbf{o}_1) \\ \mathbf{z}_1 \end{bmatrix}, \quad \text{where } \mathbf{z}_1 = T_{01}(1:3,3), \, \mathbf{o}_1 = T_{01}(1:3,4) \\
  \text{Column 3 (Joint 3):} \quad & J_3 = \begin{bmatrix} \mathbf{z}_2 \times (\mathbf{o}_{end} - \mathbf{o}_2) \\ \mathbf{z}_2 \end{bmatrix}, \quad \text{where } \mathbf{z}_2 = T_{02}(1:3,3), \, \mathbf{o}_2 = T_{02}(1:3,4) \\
  \text{Column 4 (Joint 4):} \quad & J_4 = \begin{bmatrix} \mathbf{z}_3 \times (\mathbf{o}_{end} - \mathbf{o}_3) \\ \mathbf{z}_3 \end{bmatrix}, \quad \text{where } \mathbf{z}_3 = T_{03}(1:3,3), \, \mathbf{o}_3 = T_{03}(1:3,4) \\
  \text{Column 5 (Joint 5):} \quad & J_5 = \begin{bmatrix} \mathbf{z}_4 \times (\mathbf{o}_{end} - \mathbf{o}_4) \\ \mathbf{z}_4 \end{bmatrix}, \quad \text{where } \mathbf{z}_4 = T_{04}(1:3,3), \, \mathbf{o}_4 = T_{04}(1:3,4) \\
  \text{Column 6 (Joint 6):} \quad & J_6 = \begin{bmatrix} \mathbf{z}_5 \times (\mathbf{o}_{end} - \mathbf{o}_5) \\ \mathbf{z}_5 \end{bmatrix}, \quad \text{where } \mathbf{z}_5 = T_{05}(1:3,3), \, \mathbf{o}_5 = T_{05}(1:3,4)
\end{align}

The combined configuration formula for the matrix is expressed as:
\begin{equation}
  \boxed{\;
  J(\mathbf{q}) = \begin{bmatrix} 
    \mathbf{z}_0 \times (\mathbf{o}_{end} - \mathbf{o}_0) & \dots & \mathbf{z}_5 \times (\mathbf{o}_{end} - \mathbf{o}_5) \\
    \mathbf{z}_0 & \dots & \mathbf{z}_5
  \end{bmatrix}
  \;}
\end{equation}

\end{document}

